% ==============================================================================
% Paper 86 (DE): Der Rang elliptischer Kurven: Beschränktheit und Verteilung
% Autor: Michael Fuhrmann
% Projekt: Specialist Mathematics Project, Build 167
% Datum: 2026-03-12
% ==============================================================================

\documentclass[12pt,a4paper]{amsart}

\usepackage{amsmath,amssymb,amsthm,mathtools,enumitem,booktabs,array,hyperref}
\usepackage[utf8]{inputenc}
\usepackage[T1]{fontenc}
\usepackage[ngerman]{babel}
\usepackage{geometry}
\geometry{margin=2.5cm}

% --------------------------------------------------------------------------
% Theorem-Umgebungen (Deutsch)
% --------------------------------------------------------------------------
\newtheorem{theorem}{Theorem}[section]
\newtheorem{lemma}[theorem]{Lemma}
\newtheorem{corollary}[theorem]{Korollar}
\newtheorem{proposition}[theorem]{Proposition}
\newtheorem{satz}[theorem]{Satz}
\newtheorem{korollar}[theorem]{Korollar}
\theoremstyle{definition}
\newtheorem{definition}[theorem]{Definition}
\newtheorem{definition_de}[theorem]{Definition}
\newtheorem{example}[theorem]{Beispiel}
\newtheorem{algorithmus}[theorem]{Algorithmus}
\theoremstyle{remark}
\newtheorem{remark}[theorem]{Bemerkung}
\newtheorem{bemerkung}[theorem]{Bemerkung}
\newtheorem{conjecture}[theorem]{Vermutung}
\newtheorem{vermutung}[theorem]{Vermutung}

% --------------------------------------------------------------------------
% Makros
% --------------------------------------------------------------------------
\newcommand{\Z}{\mathbb{Z}}
\newcommand{\Q}{\mathbb{Q}}
\newcommand{\R}{\mathbb{R}}
\newcommand{\C}{\mathbb{C}}
\newcommand{\F}{\mathbb{F}}
\newcommand{\N}{\mathbb{N}}
\newcommand{\Qp}{\mathbb{Q}_p}
\newcommand{\Gal}{\mathrm{Gal}}
\newcommand{\Sel}{\mathrm{Sel}}
\DeclareMathOperator{\Sha}{\mathfrak{Sh}}
\newcommand{\Prob}{\mathrm{Prob}}
\newcommand{\rank}{\mathrm{rank}}
\newcommand{\ord}{\mathrm{ord}}
\DeclareMathOperator{\Nm}{Nm}

\hypersetup{
  colorlinks=true,
  linkcolor=blue,
  citecolor=blue,
  urlcolor=blue,
  pdftitle={Der Rang elliptischer Kurven: Beschränktheit und Verteilung},
  pdfauthor={Michael Fuhrmann}
}

% ==============================================================================
\begin{document}
% ==============================================================================

\title{Der Rang elliptischer Kurven:\\
  Beschränktheit und Verteilung}

\author{Michael Fuhrmann}
\address{Specialist Mathematics Project, Build 167}
\email{specialist-maths@project.local}
\date{2026-03-12}

\begin{abstract}
Der Satz von Mordell-Weil besagt, dass für eine elliptische Kurve $E/\Q$ die Gruppe
$E(\Q)$ der rationalen Punkte endlich erzeugt ist:
$E(\Q) \cong \Z^r \oplus E(\Q)_{\mathrm{tors}}$.
Der \emph{Rang} $r = r(E)$ ist eine zentrale und rätselhafte Invariante der
Zahlentheorie.
Wir besprechen die Theorie des Rangs systematisch: den Satz von Mordell-Weil (via
Abstieg), die Selmer-Gruppe, die Tate-Shafarevich-Gruppe $\Sha(E/\Q)$, den analytischen
Rang über die $L$-Funktion $L(E,s)$ und die Birch-Swinnerton-Dyer-Vermutung (BSD),
aktuelle Rekorde (Elkies' Kurve mit Rang $\geq 29$), die Goldfeld-Vermutung
(mittlerer Rang $1/2$), die Beschränktheitsver\-mutung, Bhargava-Shankars Satz
(mittlerer Rang $\leq 0{,}885$, positive Anteile für Rang $0$ und Rang $1$),
das Poonen-Rains-Modell sowie die Paritätsvermutung.
Die Beschränktheitsvermutung und Goldfeldsche Vermutung sind beide offen.
\end{abstract}

\maketitle

\tableofcontents

% ==============================================================================
\section{Einleitung}
% ==============================================================================

Eine \emph{elliptische Kurve} über $\Q$ ist eine glatte projektive Kurve vom Geschlecht
$1$ mit einem ausgezeichneten rationalen Punkt $O$ (dem ``Punkt im Unendlichen''),
gegeben durch eine Weierstraß-Gleichung
\[
  E\colon y^2 + a_1 xy + a_3 y = x^3 + a_2 x^2 + a_4 x + a_6, \quad a_i \in \Z,
\]
mit nicht-verschwindender Diskriminante.
Die Menge $E(\Q)$ der rationalen Punkte bildet eine abelsche Gruppe
(via das Sekanten-Tangenten-Gesetz).

\begin{satz}[Mordell 1922, Weil 1928 \cite{Mordell1922,Weil1928}]
  Für jede elliptische Kurve $E/\Q$ ist die Gruppe $E(\Q)$ endlich erzeugt:
  \[
    E(\Q) \;\cong\; \Z^r \oplus E(\Q)_{\mathrm{tors}},
  \]
  wobei $r = r(E) \geq 0$ der \emph{Rang} von $E$ und $E(\Q)_{\mathrm{tors}}$ eine
  endliche abelsche Gruppe ist.
\end{satz}

Der Rang $r(E)$ ist das zentrale Objekt dieses Aufsatzes. Wesentliche Fragen:
\begin{itemize}
  \item Wie groß kann $r(E)$ werden? (Beschränktheit oder Unbeschränktheit)
  \item Was ist der ``typische'' Rang? (Verteilung, mittlerer Rang)
  \item Gibt es einen effektiven Algorithmus zur Berechnung von $r(E)$?
\end{itemize}

% ==============================================================================
\section{Der Satz von Mordell-Weil: Beweisidee via Abstieg}\label{sec:mordell_weil}
% ==============================================================================

\subsection{Schwacher Satz von Mordell-Weil}

\begin{satz}[Schwacher Satz von Mordell-Weil]
  Für jedes $n \geq 2$ ist der Quotient $E(\Q)/nE(\Q)$ endlich.
\end{satz}

\begin{proof}[Beweisskizze f\"ur $n=2$]
  Die $2$-Isogenie-Sequenz liefert eine Abbildung
  $E(\Q) \to H^1(\Gal(\overline{\Q}/\Q), E[2])$,
  die in eine endliche Selmer-Gruppe landet (durch Endlichkeit von $S$-Einheiten und
  Kummer-Theorie).
\end{proof}

\subsection{Höhe und Abstieg}

\begin{definition_de}[Kanonische Höhe (Néron-Tate)]
  Die \emph{kanonische Höhe} $\hat{h}\colon E(\Q) \to \R_{\geq 0}$ ist eine positiv
  semidefinite quadratische Form auf $E(\Q) \otimes \R$ mit $\hat{h}(P) = 0$ genau dann,
  wenn $P$ ein Torsionspunkt ist.
\end{definition_de}

Der Satz von Mordell-Weil folgt: Endliche Erzeugung ergibt sich aus dem schwachen Satz
und der Endlichkeit der Menge $\{P \in E(\Q) : \hat{h}(P) \leq C\}$ für jedes $C > 0$.

% ==============================================================================
\section{Torsionsuntergruppe}\label{sec:torsion}
% ==============================================================================

\begin{satz}[Mazur, 1977 \cite{Mazur1977}]
  Für jede elliptische Kurve $E/\Q$ ist $E(\Q)_{\mathrm{tors}}$ isomorph zu einer der
  $15$ Gruppen:
  \begin{align*}
    &\Z/n\Z \quad (n = 1,2,\ldots,10,12), \\
    &\Z/2\Z \times \Z/2n\Z \quad (n = 1,2,3,4).
  \end{align*}
\end{satz}

\begin{bemerkung}
  Jede dieser $15$ Torsionsstrukturen wird von einer konkreten elliptischen Kurve über
  $\Q$ realisiert. Die häufigsten sind $\Z/1\Z$ (trivial) und $\Z/2\Z$.
\end{bemerkung}

% ==============================================================================
\section{Selmer-Gruppe und Tate-Shafarevich-Gruppe}\label{sec:selmer}
% ==============================================================================

\begin{definition_de}[Selmer-Gruppe und $\Sha$]
  Für eine Primzahl $p$ liefert die kurze exakte Sequenz
  $0 \to E[p] \to E \xrightarrow{[p]} E \to 0$
  via Galois-Kohomologie:
  \[
    0 \;\to\; \frac{E(\Q)}{pE(\Q)} \;\to\; \Sel^{(p)}(E/\Q)
    \;\to\; \Sha(E/\Q)[p] \;\to\; 0.
  \]
  Die \emph{$p$-Selmer-Gruppe} ist
  \[
    \Sel^{(p)}(E/\Q) = \ker\!\left(H^1(\Q, E[p]) \to \prod_v H^1(\Q_v, E)\right).
  \]
  Die \emph{Tate-Shafarevich-Gruppe} ist
  \[
    \Sha(E/\Q) = \ker\!\left(H^1(\Q, E) \to \prod_v H^1(\Q_v, E)\right).
  \]
\end{definition_de}

\begin{vermutung}[Endlichkeit von $\Sha$]\label{verm:sha}
  Für jede elliptische Kurve $E/\Q$ ist $\Sha(E/\Q)$ endlich.
\end{vermutung}

\begin{bemerkung}
  Endlichkeit von $\Sha$ ist nur in Spezialfällen bewiesen:
  \begin{itemize}
    \item wenn $\ord_{s=1}L(E,s) = 0$: Kolyvagin \cite{Kolyvagin1988} bewies $\rank E(\Q) = 0$ und $\Sha(E/\Q)$ endlich (mittels Euler-Systemen);
    \item wenn $\ord_{s=1}L(E,s) = 1$: Gross-Zagier \cite{GZ1986} + Kolyvagin bewiesen $\rank E(\Q) = 1$ und $\Sha(E/\Q)$ endlich;
    \item für spezifische Familien via Heegner-Punkte.
  \end{itemize}
  Im Allgemeinen ist Vermutung~\ref{verm:sha} offen.
\end{bemerkung}

% ==============================================================================
\section{$L$-Funktion und die BSD-Vermutung}\label{sec:bsd}
% ==============================================================================

\begin{definition_de}[$L$-Funktion einer elliptischen Kurve]
  Für eine Primzahl $p$ guter Reduktion ist der lokale $L$-Faktor
  \[
    L_p(E,s) = \frac{1}{1 - a_p p^{-s} + p^{1-2s}}, \quad
    a_p = p + 1 - \#E(\F_p).
  \]
  Die globale $L$-Funktion ist $L(E,s) = \prod_p L_p(E,s)$,
  konvergent für $\mathrm{Re}(s) > 3/2$, analytisch fortsetzbar auf ganz $\C$
  (durch Modularität, Wiles et al.\ \cite{BCDT2001}).
\end{definition_de}

\begin{vermutung}[Birch und Swinnerton-Dyer \cite{BSD1965}]\label{verm:bsd}
  Sei $E/\Q$ eine elliptische Kurve. Dann:
  \begin{enumerate}[label=(\roman*)]
    \item $\ord_{s=1} L(E,s) = \rank E(\Q)$ \quad (analytischer Rang $=$ algebraischer Rang);
    \item Der Leitkoeffizient von $L(E,s)$ bei $s=1$ ist durch die BSD-Formel gegeben:
    \[
      \lim_{s\to 1} \frac{L(E,s)}{(s-1)^r} =
      \frac{\Omega_E \cdot \mathrm{Reg}(E/\Q) \cdot |\Sha(E/\Q)| \cdot \prod_p c_p}
           {|E(\Q)_{\mathrm{tors}}|^2},
    \]
    wobei $\Omega_E$ die reelle Periode, $\mathrm{Reg}(E/\Q)$ der Néron-Tate-Regulator,
    $c_p$ die Tamagawa-Zahlen und $r = \rank E(\Q)$.
  \end{enumerate}
\end{vermutung}

\begin{bemerkung}
  BSD ist eines der Millennium-Preisaufgaben des Clay Mathematics Institute.
  Das beste bekannte Ergebnis ist:
\end{bemerkung}

\begin{satz}[Coates-Wiles \cite{CW1977}, Kolyvagin \cite{Kolyvagin1988},
  Gross-Zagier \cite{GZ1986}]
  Gilt $L(E,1) \neq 0$, so ist $\rank E(\Q) = 0$.
  Gilt $L(E,1) = 0$ und $L'(E,1) \neq 0$ (analytischer Rang $1$) und $E$ hat einen
  Heegner-Punkt unendlicher Ordnung, so ist $\rank E(\Q) = 1$ und $\Sha(E/\Q)$ endlich.
\end{satz}

% ==============================================================================
\section{Rekorde: Elliptische Kurven mit hohem Rang}\label{sec:rekorde}
% ==============================================================================

\begin{satz}[Elkies, 2006 \cite{Elkies2006}]
  Es gibt eine elliptische Kurve über $\Q$ mit Rang $\geq 29$.
  Der aktuelle Rekord (Stand 2026) ist Rang $\geq 29$, gefunden von Noam Elkies.
\end{satz}

\begin{bemerkung}
  Frühere Rekorde: Mestre \cite{Mestre1982} (Rang $\geq 12$, 1982),
  Fermigier (Rang $\geq 19$, 1997),
  Martin-McMillen (Rang $\geq 24$, 2000),
  Elkies (Rang $\geq 28$ und $\geq 29$, 2006).
\end{bemerkung}

% ==============================================================================
\section{Die Goldfeld-Vermutung und die Beschränktheits-Vermutung}\label{sec:goldfeld}
% ==============================================================================

\begin{vermutung}[Goldfeld \cite{Goldfeld1979}]\label{verm:goldfeld}
  Für elliptische Kurven $E/\Q$, geordnet nach dem Führer $N$, gilt:
  \[
    \lim_{X\to\infty} \frac{\sum_{N(E) \leq X} \rank E(\Q)}
                           {\#\{E : N(E) \leq X\}}
    \;=\; \frac{1}{2}.
  \]
  Äquivalent: $50\,\%$ der elliptischen Kurven haben Rang $0$, $50\,\%$ haben Rang $1$;
  der Anteil der Kurven mit Rang $\geq 2$ ist null.
\end{vermutung}

\begin{bemerkung}
  Der Erwartungswert $1/2$ folgt aus der Paritätsvermutung: Näherungsweise die Hälfte
  der Kurven hat geraden analytischen Rang (Vorzeichen $w(E) = +1$) und die andere Hälfte
  ungeraden ($w(E) = -1$).
\end{bemerkung}

\begin{vermutung}[Beschränktheit des Rangs]\label{verm:beschraenkt}
  Es gibt eine absolute Konstante $C > 0$, so dass für alle elliptischen Kurven $E/\Q$
  gilt: $\rank E(\Q) \leq C$.
\end{vermutung}

\begin{bemerkung}
  Diese Vermutung ist sehr umstritten: Heuristiken von Poonen-Rains \cite{PR2012} und
  Park-Poonen-Voight-Wood \cite{PPVW2019} sprechen dafür; es ist aber kein theoretisches
  Hindernis für beliebig großen Rang bekannt.
\end{bemerkung}

% ==============================================================================
\section{Bhargava-Shankar: Schranken für den mittleren Rang}\label{sec:bhargava}
% ==============================================================================

\begin{satz}[Bhargava-Shankar \cite{BS2015}]\label{satz:bs}
  Werden elliptische Kurven $E/\Q$ nach der Höhe $H(E) = \max(|a|^3, |b|^2)$
  für das Modell $y^2 = x^3 + ax + b$ geordnet, so gilt:
  \begin{enumerate}[label=(\roman*)]
    \item Der mittlere Umfang der $2$-Selmer-Gruppe $\Sel^{(2)}(E/\Q)$ ist $3$.
    \item Der mittlere Rang erfüllt $\overline{r} \leq 0{,}885$.
    \item Ein positiver Anteil der elliptischen Kurven hat Rang $0$.
    \item Ein positiver Anteil der elliptischen Kurven hat Rang $1$.
  \end{enumerate}
\end{satz}

\begin{bemerkung}
  Die Schranke $\overline{r} \leq 0{,}885$ ergibt sich aus dem mittleren Umfang der
  $2$-Selmer-Gruppe ($= 3$) und des $3$-Selmer-Gruppe ($= 4$) kombiniert mit dem
  $5$-Abstieg.
\end{bemerkung}

\begin{satz}[Bhargava-Skinner-Zhang \cite{BSZ2014}]
  Ein positiver Anteil der elliptischen Kurven hat $\rank E(\Q) = 0$ und
  $|\Sha(E/\Q)| \mid 4$.
  Ein positiver Anteil hat $\rank E(\Q) = 1$ mit endlichem $|\Sha|$
  (bedingt unter BSD).
\end{satz}

% ==============================================================================
\section{Die Paritätsvermutung}\label{sec:paritaet}
% ==============================================================================

\begin{definition_de}[Vorzeichen / Rootnumber]
  Das \emph{Vorzeichen} $w(E) \in \{+1, -1\}$ tritt in der Funktionalgleichung von
  $L(E,s)$ auf:
  \[
    \Lambda(E,s) = w(E) \cdot \Lambda(E, 2-s),
  \]
  wo $\Lambda(E,s) = N^{s/2}(2\pi)^{-s}\Gamma(s) L(E,s)$ die vervollständigte
  $L$-Funktion ist.
\end{definition_de}

\begin{vermutung}[Paritätsvermutung]\label{verm:paritaet}
  Für jede elliptische Kurve $E/\Q$ gilt:
  \[
    (-1)^{\rank E(\Q)} \;=\; w(E).
  \]
\end{vermutung}

\begin{satz}[Nekovář \cite{Nekovar2009}]
  Die Paritätsvermutung gilt für elliptische Kurven über total-reellen Körpern $F$,
  für die $E$ bei allen Primidealen über $2$ potenziell semistabil ist
  (unter gewissen lokalen Bedingungen).
\end{satz}

% ==============================================================================
\section{Das Poonen-Rains-Modell}\label{sec:poonen_rains}
% ==============================================================================

\begin{vermutung}[Poonen-Rains-Modell \cite{PR2012}]\label{verm:pr}
  Sei $p$ prim. Die ``Wahrscheinlichkeit'', dass eine ``zufällige'' elliptische Kurve
  $E/\Q$ $p$-Selmer-Rang $\geq r$ hat, wird durch das Zufallsmatrixmodell beschrieben:
  \[
    \Prob[\dim_{\F_p} \Sel^{(p)}(E/\Q) \geq r] \;\sim\;
    \prod_{k=1}^\infty (1 + p^{-k})^{-1} \cdot p^{-r(r-1)/2},
  \]
  moduliert durch das Vorzeichen.
\end{vermutung}

\begin{bemerkung}
  Unter den Poonen-Rains-Heuristiken gilt:
  \begin{itemize}
    \item Der mittlere Rang ist $1/2$ (konsistent mit Goldfeldsche Vermutung);
    \item $50\,\%$ der Kurven haben Rang $0$, $50\,\%$ Rang $1$;
    \item Der Anteil der Rang-$\geq 2$-Kurven ist null (dichte null, aber unendlich viele);
    \item Der Anteil der Rang-$\geq r$-Kurven fällt super-exponentiell: $\sim p^{-r^2}$.
  \end{itemize}
\end{bemerkung}

% ==============================================================================
\section{Brumer-McGuinness-Daten}\label{sec:daten}
% ==============================================================================

\begin{bemerkung}
  Brumer und McGuinness \cite{BM1990} führten eine der ersten groß angelegten
  rechnerischen Untersuchungen des Rangs durch. Für elliptische Kurven
  $y^2 = x^3 + ax + b$ mit $|a|, |b| \leq B$ ergaben sich ungefähr:
  \begin{center}
    \begin{tabular}{lr}
      \toprule
      Rang & Anteil \\
      \midrule
      0 & $\approx 35{,}2\,\%$ \\
      1 & $\approx 42{,}7\,\%$ \\
      2 & $\approx 18{,}9\,\%$ \\
      $\geq 3$ & $\approx 3{,}2\,\%$ \\
      \bottomrule
    \end{tabular}
  \end{center}
  Diese Daten sind konsistent mit der Goldfeldsche Vermutung (für ausreichend große $B$
  sollten die Anteile von Rang $0$ und $1$ gegen $50\,\%$ konvergieren).
\end{bemerkung}

% ==============================================================================
\section{2-Abstieg und Algorithmen}\label{sec:abstieg}
% ==============================================================================

\begin{definition_de}[$2$-Abstieg]
  Ein \emph{$2$-Abstieg} auf $E/\Q$ berechnet die $2$-Selmer-Gruppe
  $\Sel^{(2)}(E/\Q)$ explizit.
  Gegeben $E\colon y^2 = (x-e_1)(x-e_2)(x-e_3)$ (mit $e_i \in \Q$ paarweise
  verschieden), entsprechen die $2$-Selmer-Elemente quadratischen Verdrillungen
  $E_d$ mit rationalen Punkten (homogene Räume).
\end{definition_de}

\begin{algorithmus}[$2$-Abstiegs-Algorithmus]
  \begin{enumerate}[label=(\arabic*)]
    \item Berechne $E[2](\Q)$ (die $2$-Torsion).
    \item Für jede Klasse $[c] \in (\Q^\times / (\Q^\times)^2)^2$, eingeschränkt auf
      $S$-Einheiten ($S$ = Primzahlen schlechter Reduktion), prüfe, ob der entsprechende
      homogene Raum einen rationalen Punkt hat.
    \item Gib $\Sel^{(2)}(E/\Q)$ aus als die Klassen mit überall lokalen Punkten.
  \end{enumerate}
\end{algorithmus}

\begin{bemerkung}
  Der $2$-Selmer-Rang ist eine obere Schranke für $\rank E(\Q) + \dim_{\F_2} \Sha[2]$.
  Um $\rank E(\Q)$ exakt zu bestimmen, muss man $\Sha[2] = 0$ nachweisen
  (was nicht immer beweisbar ist).
\end{bemerkung}

% ==============================================================================
\section{Offene Probleme}\label{sec:offene_probleme}
% ==============================================================================

\begin{vermutung}[BSD, volle Version; Vermutung~\ref{verm:bsd}]
  Beide Teile von BSD sind offen.
  Nur für analytischen Rang $\leq 1$ sind Teilresultate bekannt.
\end{vermutung}

\begin{vermutung}[Goldfeld-Vermutung; Vermutung~\ref{verm:goldfeld}]
  Mittlerer Rang $= 1/2$.
  Bestes bekanntes Ergebnis: $\leq 0{,}885$ (Bhargava-Shankar).
\end{vermutung}

\begin{vermutung}[Beschränktheitsvermutung; Vermutung~\ref{verm:beschraenkt}]
  Ist der Rang uniform beschränkt?
  Aktueller Rekord: $\geq 29$ (Elkies).
  Kein theoretisches Hindernis für beliebig großen Rang bekannt.
\end{vermutung}

\begin{vermutung}[Endlichkeit von $\Sha$; Vermutung~\ref{verm:sha}]
  Im Allgemeinen offen; nur für Rang $\leq 1$ bedingt bekannt.
\end{vermutung}

% ==============================================================================
\section{Fazit}\label{sec:fazit}
% ==============================================================================

Der Rang elliptischer Kurven ist eines der meistuntersuchten und rätselhaftesten
Invariante der Zahlentheorie. Der Satz von Mordell-Weil sichert seine Endlichkeit,
aber die Berechnung von $r(E)$ ist im Allgemeinen ein schwieriges offenes Problem
(kein unbedingter Algorithmus bekannt). Die BSD-Vermutung, falls wahr, würde einen
analytischen Zugang zu $r(E)$ liefern. Die Arbeit von Bhargava und Shankar gibt die
erste statistische Kontrolle über die Rangverteilung: Der mittlere Rang ist höchstens
$0{,}885$, und positive Anteile von Rang-$0$- und Rang-$1$-Kurven wurden nachgewiesen.
Die Fragen nach Beschränktheit des Rangs und nach dem genauen mittleren Rang
bleiben zu den tiefsten ungelösten Problemen der arithmetischen Geometrie.

\begin{thebibliography}{99}

\bibitem{BS2015}
M.\ Bhargava, A.\ Shankar,
\textit{Binary quartic forms having bounded invariants, and the boundedness of the
average rank of elliptic curves},
Ann.\ of Math.\ (2) \textbf{181} (2015), Nr.\ 1, 191--242.

\bibitem{BS2015a}
M.\ Bhargava, A.\ Shankar,
\textit{Ternary cubic forms having bounded invariants, and the existence of a positive
proportion of elliptic curves having rank $0$},
Ann.\ of Math.\ (2) \textbf{181} (2015), Nr.\ 2, 587--621.

\bibitem{BSZ2014}
M.\ Bhargava, C.\ Skinner, W.\ Zhang,
\textit{A majority of elliptic curves over $\mathbb{Q}$ satisfy the Birch and
Swinnerton-Dyer conjecture},
Preprint (2014), arXiv:1407.1826.

\bibitem{BM1990}
A.\ Brumer, O.\ McGuinness,
\textit{The behavior of the Mordell-Weil group of elliptic curves},
Bull.\ Amer.\ Math.\ Soc.\ \textbf{23} (1990), Nr.\ 2, 375--382.

\bibitem{BSD1965}
B.\,J.\ Birch, H.\,P.\,F.\ Swinnerton-Dyer,
\textit{Notes on elliptic curves. II},
J.\ Reine Angew.\ Math.\ \textbf{218} (1965), 79--108.

\bibitem{BCDT2001}
C.\ Breuil, B.\ Conrad, F.\ Diamond, R.\ Taylor,
\textit{On the modularity of elliptic curves over $\mathbf{Q}$: wild $3$-adic exercises},
J.\ Amer.\ Math.\ Soc.\ \textbf{14} (2001), Nr.\ 4, 843--939.

\bibitem{CW1977}
J.\ Coates, A.\ Wiles,
\textit{On the conjecture of Birch and Swinnerton-Dyer},
Invent.\ Math.\ \textbf{39} (1977), Nr.\ 3, 223--251.

\bibitem{Elkies2006}
N.\ Elkies,
\textit{$\mathbb{Z}^{28}$ in $E(\mathbb{Q})$},
Eintrag in der NMBRTHRY-Liste (Mai 2006).

\bibitem{Goldfeld1979}
D.\ Goldfeld,
\textit{Conjectures on elliptic curves over quadratic fields},
in \textit{Number Theory, Carbondale 1979}, Lect.\ Notes Math.\ \textbf{751},
Springer (1979), 108--118.

\bibitem{GZ1986}
B.\,H.\ Gross, D.\,B.\ Zagier,
\textit{Heegner points and derivatives of $L$-series},
Invent.\ Math.\ \textbf{84} (1986), Nr.\ 2, 225--320.

\bibitem{Kolyvagin1988}
V.\,A.\ Kolyvagin,
\textit{Finiteness of $E(\mathbb{Q})$ and $\Sha(E,\mathbb{Q})$ for a subclass of
Weil curves},
Izv.\ Akad.\ Nauk SSSR Ser.\ Mat.\ \textbf{52} (1988), 522--540.

\bibitem{Mazur1977}
B.\ Mazur,
\textit{Modular curves and the Eisenstein ideal},
Publ.\ Math.\ IHÉS \textbf{47} (1977), 33--186.

\bibitem{Mestre1982}
J.-F.\ Mestre,
\textit{Construction d'une courbe elliptique de rang $\geq 12$},
C.\ R.\ Acad.\ Sci.\ Paris \textbf{295} (1982), 643--644.

\bibitem{Mordell1922}
L.\,J.\ Mordell,
\textit{On the rational solutions of the indeterminate equation of the third and fourth degrees},
Proc.\ Cambridge Philos.\ Soc.\ \textbf{21} (1922), 179--192.

\bibitem{Nekovar2009}
J.\ Nekovář,
\textit{On the parity of ranks of Selmer groups IV},
Compos.\ Math.\ \textbf{145} (2009), Nr.\ 6, 1351--1359.

\bibitem{PPVW2019}
J.\ Park, B.\ Poonen, J.\ Voight, M.\,M.\ Wood,
\textit{A heuristic for boundedness of ranks of elliptic curves},
J.\ Eur.\ Math.\ Soc.\ \textbf{21} (2019), Nr.\ 9, 2859--2903.

\bibitem{PR2012}
B.\ Poonen, E.\ Rains,
\textit{Random maximal isotropic subspaces and Selmer groups},
J.\ Amer.\ Math.\ Soc.\ \textbf{25} (2012), Nr.\ 1, 245--269.

\bibitem{Weil1928}
A.\ Weil,
\textit{L'arithmétique sur les courbes algébriques},
Acta Math.\ \textbf{52} (1928), 281--315.

\end{thebibliography}

\end{document}
